%
% Capítulo 9
%
\chapter{Conclusões} \label{cap:conclusoes}

O desenvolvimento do projeto \textit{Sistema de Planeamento de Horários} permitiu demonstrar a viabilidade técnica da solução proposta e validar com sucesso todos os seus fluxos funcionais. Ao longo do projeto foi possível consolidar uma arquitetura robusta composta por aplicação móvel, backend e persistência remota, inserindo a solução num contexto de utilização real.

A solução final entregue suporta na íntegra as operações centrais planeadas, como o registo e autenticação de utilizadores, a associação entre alunos e professores, o registo de disponibilidades, a definição de restrições complexas e a construção automatizada de propostas de horário. A integração bem-sucedida da aceitação do horário com o \textit{Google Calendar} do professor acrescentou um valor prático imediato à solução, reforçando a sua utilidade no quotidiano de um centro de estudos.

Para além da implementação funcional, o projeto permitiu validar com rigor a comunicação síncrona entre frontend e backend, otimizar a usabilidade da interface móvel e consolidar mecanismos críticos de segurança, incluindo a autenticação remota e o armazenamento seguro de credenciais. Os testes de integração automatizados realizados demonstraram que a solução responde de forma coerente e estável a todos os casos de uso identificados.

O trabalho desenvolvido permitiu ainda abrir uma linha exploratória complementar inovadora relacionada com a interpretação de disponibilidades a partir de inputs não estruturados, com recurso a modelos multimodais de Inteligência Artificial. Embora esta componente não tenha sido integrada na aplicação final devido à priorização da estabilidade do sistema base, o estudo exploratório efetuado permitiu identificar abordagens robustas e documentadas para o tratamento de texto, imagem e áudio em futuras evoluções comerciais do produto.

Em síntese, o produto final demonstra utilidade prática real, coerência arquitetural e cumpre integralmente os objetivos pedagógicos e técnicos delineados nos objetivos deste projeto.
